\documentclass[11pt]{article}
\usepackage[margin=1in]{geometry}
\usepackage{amsmath}
\usepackage{booktabs}
\usepackage{hyperref}
\usepackage{url}

\title{Memo 09: Spatial Distribution Analysis for FVCOM-MP Cases C and D\\
       Effective Floc Coupling, Days 10 and 30}
\author{FVCOM-MP WaterPACT Investigation}
\date{May 2026}

\begin{document}
\sloppy
\maketitle

\section*{Purpose}

This memo records the first spatial-distribution comparison between case
\texttt{c}, the comparison MP-floc case, and case \texttt{d}, the
ORIG\_SED effective-floc coupling case.  The analysis used the day-10 and
day-30 MAT files already extracted in \texttt{OUTPUT\_c} and
\texttt{OUTPUT\_d}.  Day 90 is not included here because the corresponding
v2 MAT files were not present at the time of this memo.

\section*{Files and Command}

The analysis inputs were:
\begin{itemize}
  \item \texttt{FVCOM\_MP\_test\_run/OUTPUT\_c/spatial\_distribution\_v2\_c\_day10.mat}
  \item \texttt{FVCOM\_MP\_test\_run/OUTPUT\_d/spatial\_distribution\_v2\_d\_day10.mat}
  \item \texttt{FVCOM\_MP\_test\_run/OUTPUT\_c/spatial\_distribution\_v2\_c\_day30.mat}
  \item \texttt{FVCOM\_MP\_test\_run/OUTPUT\_d/spatial\_distribution\_v2\_d\_day30.mat}
\end{itemize}

The analysis was run with:
\begin{verbatim}
python FVCOM_MP_test_run/PYTHON/analyze_spatial_distribution_v2.py --days 10 30
\end{verbatim}

The outputs were written to:
\begin{itemize}
  \item \texttt{FVCOM\_MP\_test\_run/PYTHON/output/spatial\_distribution\_v2/day10}
  \item \texttt{FVCOM\_MP\_test\_run/PYTHON/output/spatial\_distribution\_v2/day30}
\end{itemize}
Each day folder contains per-case maps, \texttt{d-c} maps,
\texttt{summary\_stats.csv}, \texttt{diff\_summary\_stats.csv}, and
\texttt{README.md}.  The run generated 108 PNG maps across the two day folders.

\section*{Extraction Window}

Both day-10 and day-30 MAT files contain 24 output records.  The day-10 files
span simulation day 9.0000 to 9.9570, and the day-30 files span simulation day
29.0000 to 29.9570.  These are therefore final-24-hour averages before the
target day endpoints, as intended for this v2 spatial workflow.

\section*{Equivalent Floc Size Post-processing}

The equivalent floc size was computed from the existing output
\texttt{settle\_vel\_floc\_mp1} with
\begin{equation}
  d_{\mathrm{floc,eff}} =
  \left[
  \frac{18 \mu W_{\mathrm{floc}}}
       {g(\rho_{\mathrm{floc,eff}}-\rho_w)}
  \right]^{1/2},
\end{equation}
where $\rho_{\mathrm{floc,eff}} = 1300$ kg m$^{-3}$ and
\begin{equation}
  \nu = 10^{-6}\exp[-0.025(T-20)], \qquad
  0.5\times10^{-6} \le \nu \le 2.0\times10^{-6},
\end{equation}
with $\mu=\rho_w\nu$.  This is the same post-processing relation documented in
Memo 08.

\section*{Key Spatial Statistics}

Table~\ref{tab:mainstats} summarizes the domain means and 95th percentiles for
the depth-averaged fields most relevant to the effective-floc coupling.  All
values are taken directly from the generated \texttt{summary\_stats.csv} files.

\begin{table}[h]
\centering
\caption{Depth-averaged domain statistics for selected fields.}
\label{tab:mainstats}
\begin{tabular}{llrrrr}
\toprule
Day & Field & C mean & D mean & C p95 & D p95 \\
\midrule
10 & Sediment & $1.035$ & $4.506$ & $4.332$ & $10.506$ \\
10 & $W_{\mathrm{floc}}$ (m s$^{-1}$) & $1.463{\times}10^{-3}$ & $7.527{\times}10^{-5}$ & $3.000{\times}10^{-3}$ & $2.079{\times}10^{-4}$ \\
10 & $d_{\mathrm{floc,eff}}$ ($\mu$m) & $79.07$ & $19.35$ & $163.17$ & $38.69$ \\
10 & $\lambda_a$ (s$^{-1}$) & $2.110{\times}10^{-2}$ & $0$ & $8.713{\times}10^{-2}$ & $0$ \\
10 & $\lambda_d$ (s$^{-1}$) & $1.943{\times}10^{-2}$ & $4.996{\times}10^{-1}$ & $6.999{\times}10^{-2}$ & $5.000{\times}10^{-1}$ \\
10 & MP total (kg m$^{-3}$) & $1.162{\times}10^{-6}$ & $6.697{\times}10^{-7}$ & $4.344{\times}10^{-6}$ & $2.700{\times}10^{-6}$ \\
10 & MP aggregated fraction & $0.363$ & $0$ & $0.433$ & $0$ \\
\midrule
30 & Sediment & $0.524$ & $4.373$ & $2.518$ & $10.143$ \\
30 & $W_{\mathrm{floc}}$ (m s$^{-1}$) & $2.720{\times}10^{-3}$ & $1.174{\times}10^{-4}$ & $3.000{\times}10^{-3}$ & $2.809{\times}10^{-4}$ \\
30 & $d_{\mathrm{floc,eff}}$ ($\mu$m) & $148.31$ & $30.89$ & $167.08$ & $47.62$ \\
30 & $\lambda_a$ (s$^{-1}$) & $1.028{\times}10^{-2}$ & $0$ & $4.213{\times}10^{-2}$ & $0$ \\
30 & $\lambda_d$ (s$^{-1}$) & $8.798{\times}10^{-3}$ & $4.999{\times}10^{-1}$ & $3.934{\times}10^{-2}$ & $5.000{\times}10^{-1}$ \\
30 & MP total (kg m$^{-3}$) & $3.495{\times}10^{-5}$ & $1.267{\times}10^{-6}$ & $1.702{\times}10^{-4}$ & $3.721{\times}10^{-6}$ \\
30 & MP aggregated fraction & $0.181$ & $0$ & $0.388$ & $0$ \\
\bottomrule
\end{tabular}
\end{table}

\section*{Main Findings}

\subsection*{The effective-floc settling signal changed strongly and plausibly}

Case \texttt{d} has much smaller \texttt{settle\_vel\_floc\_mp1} than case
\texttt{c}.  The domain-mean depth-averaged floc settling speed decreases from
$1.46\times10^{-3}$ to $7.53\times10^{-5}$ m s$^{-1}$ on day 10, and from
$2.72\times10^{-3}$ to $1.17\times10^{-4}$ m s$^{-1}$ on day 30.  The maximum
case-\texttt{d} depth-averaged values were $4.32\times10^{-4}$ and
$4.87\times10^{-4}$ m s$^{-1}$ on days 10 and 30, respectively, which is
consistent with the ORIG\_SED floc settling cap of approximately
$5.0\times10^{-4}$ m s$^{-1}$.

\subsection*{The post-processed effective floc diameters are much smaller}

The Stokes-equivalent floc size also decreases sharply in case \texttt{d}.  The
domain-mean depth-averaged effective diameter is about 19.4 $\mu$m on day 10
and 30.9 $\mu$m on day 30.  The corresponding 95th percentiles are about
38.7 $\mu$m and 47.6 $\mu$m.  The maximum depth-averaged case-\texttt{d}
values are approximately 61.6 $\mu$m on day 10 and 66.0 $\mu$m on day 30.

\subsection*{The present size gate disables aggregation in case D}

The plastic input has \texttt{PLAST\_AAXI = PLAST\_BAXI = PLAST\_CAXI = 0.100},
which is read as a 0.100 mm, or 100 $\mu$m, equivalent plastic diameter.  In the
current aggregation kernel, the critical MP size is computed from floc size as
\begin{equation}
  d_{\mathrm{MP,max}} =
  -0.0002 d_{\mathrm{floc}}^2 + 0.36 d_{\mathrm{floc}},
  \quad d_{\mathrm{floc}} < 900\ \mu\mathrm{m}.
\end{equation}
For the case-\texttt{d} effective floc range of roughly 20--66 $\mu$m, this
gives $d_{\mathrm{MP,max}}$ of only about 7--23 $\mu$m.  The present 100 $\mu$m
MP class is therefore larger than the allowed size for aggregation.  This
explains why \texttt{lambda\_a\_mp1} is zero everywhere in case \texttt{d}, why
\texttt{mp1\_agg} is zero, and why the aggregated MP fraction is zero in both
day-10 and day-30 outputs.

\subsection*{Disaggregation rate is high, but there is no aggregated MP mass}

Case \texttt{d} has \texttt{lambda\_d\_mp1} close to 0.5 s$^{-1}$ across the
domain.  Because \texttt{mp1\_agg} is zero in case \texttt{d}, this rate should
be read mainly as a diagnostic rate produced by the kernel under the current
effective-floc sizes, not as evidence of an active disaggregation flux.

\subsection*{Sediment fields differ substantially between cases}

The case-\texttt{d} sediment concentration is higher in the depth average on
both days: the mean is $4.51$ versus $1.03$ on day 10, and $4.37$ versus
$0.524$ on day 30.  Surface sediment is broadly higher in case \texttt{d},
while bottom sediment has mixed differences and can be lower than case
\texttt{c}.  This is expected to reflect the changed cohesive ORIG\_SED
sediment configuration in case \texttt{d}, not only the MP-side effective-floc
switch.

\subsection*{MP response is strong}

Case \texttt{d} has no aggregated MP and no bottom MP mass in the analyzed
outputs.  The domain-mean total suspended MP concentration is lower than case
\texttt{c}: about 42\% lower on day 10 and about 96\% lower on day 30.  This
suggests that, under the present 100 $\mu$m MP size and ORIG-equivalent floc
sizes, the effective-floc coupling moves the model into a non-aggregation
regime rather than a modified-aggregation regime.

\section*{Interpretation}

The output confirms that the new post-processing workflow can see the intended
case-\texttt{d} ORIG\_SED settling signal through
\texttt{settle\_vel\_floc\_mp1}.  The effective settling speed and the
back-calculated equivalent floc size both change in a physically consistent
direction: lower settling speed produces smaller Stokes-equivalent flocs.

The important model-behavior result is that these smaller effective flocs are
too small to aggregate the current 100 $\mu$m plastic class under the existing
critical-size closure.  Therefore, the case-\texttt{d} run is not just changing
the aggregation strength; it is effectively shutting aggregation off for this
plastic class.

\section*{Sediment Setup Comparison and Equivalence Issue}

The case-\texttt{c} and case-\texttt{d} results above should not be interpreted
as an isolated test of only \texttt{PLAST\_FLOC\_EFFECTIVE}.  The sediment
configuration also changed between the two cases.  Table~\ref{tab:sedsetup}
summarizes the active sediment-module inputs.

\begin{table}[h]
\centering
\caption{Active sediment-module setup in cases \texttt{c} and \texttt{d}.}
\label{tab:sedsetup}
\begin{tabular}{lll}
\toprule
Parameter & Case \texttt{c} & Case \texttt{d} \\
\midrule
Sediment file &
\texttt{generic\_sediment\_c.inp} &
\texttt{generic\_sediment\_d.inp} \\
\texttt{SED\_TYPE} &
\texttt{non-cohesive} &
\texttt{cohesive} \\
\texttt{SED\_SD50} &
1.0 mm &
0.004 mm \\
\texttt{SED\_WSET} input &
3.0 mm s$^{-1}$ &
0.1 mm s$^{-1}$ \\
Runtime base settling speed &
$3.0\times10^{-3}$ m s$^{-1}$ &
$1.0\times10^{-4}$ m s$^{-1}$ \\
Settling treatment &
constant non-cohesive \texttt{Wset} &
ORIG\_SED cohesive flocculation \\
\bottomrule
\end{tabular}
\end{table}

In \texttt{mod\_sed.F}, \texttt{SED\_WSET} is read in mm s$^{-1}$ and converted
to m s$^{-1}$.  For non-cohesive sediment, the settling calculation uses
\texttt{wset = sed(ised)\%Wset}.  For cohesive sediment, the code instead calls
\texttt{Calc\_Wset\_Floc} with \texttt{sed(ised)\%w0} as the base settling
velocity and then applies the ORIG flocculation treatment.  Therefore, case
\texttt{c} uses a constant non-cohesive settling speed of 0.003 m s$^{-1}$,
whereas case \texttt{d} starts from a cohesive base settling speed of
$1.0\times10^{-4}$ m s$^{-1}$ and then reconstructs the ORIG floc signal.

This means the current \texttt{c}--\texttt{d} comparison is not dynamically
equivalent.  Case \texttt{d} simultaneously changes:
\begin{itemize}
  \item sediment type from non-cohesive to cohesive;
  \item primary sediment size from 1.0 mm to 0.004 mm;
  \item base settling speed from 3.0 mm s$^{-1}$ to 0.1 mm s$^{-1}$;
  \item the sediment settling closure from constant settling to ORIG cohesive
        flocculation;
  \item the MP interaction path through \texttt{PLAST\_FLOC\_EFFECTIVE=T}.
\end{itemize}

The plastic input files also contain a separate static parameter
\texttt{SED\_SD50 = 0.35} mm and \texttt{PLAST\_WSSED = 0.01} m s$^{-1}$ in
both cases.  These are plastic-module ambient or bed parameters and are not the
active suspended sediment class settings used by \texttt{mod\_sed.F}.  They
should not be used to judge equivalence of the active sediment transport cases.

For a clean MP-coupling comparison, the control and switch-on cases should keep
the active sediment module configuration identical.  One practical setup would
be:
\begin{enumerate}
  \item define a cohesive control case with the same \texttt{SED\_TYPE},
        \texttt{SED\_SD50}, \texttt{SED\_WSET}, erosion, deposition, bed, and
        forcing settings as case \texttt{d};
  \item keep \texttt{PLAST\_FLOC=T} but set
        \texttt{PLAST\_FLOC\_EFFECTIVE=F} in that control case;
  \item compare it against a matched case with only
        \texttt{PLAST\_FLOC\_EFFECTIVE=T}.
\end{enumerate}
That paired design would isolate the MP-side effective-floc-size treatment.
The present case \texttt{c} versus case \texttt{d} comparison is still useful as
a broad sediment-regime sensitivity, but it cannot support a strict
attribution of the MP response to the effective-floc switch alone.

\section*{Recommended Next Checks}

\begin{enumerate}
  \item Confirm that the 0.100 mm MP size is the intended particle size for
        this sensitivity.  If the intention is to test ORIG-effective floc
        aggregation, a smaller MP class may be needed.
  \item Consider a sensitivity to \texttt{PLAST\_FLOC\_RHO\_EFF}.  Lowering the
        assumed effective floc density toward water density increases the
        Stokes-equivalent diameter for the same settling speed.
  \item Consider adding optional diagnostics for
        $d_{\mathrm{floc,eff}}$ and $d_{\mathrm{MP,max}}$ directly in model
        output.  The current post-processing can recover
        $d_{\mathrm{floc,eff}}$, but $d_{\mathrm{MP,max}}$ is still inferred
        from the aggregation closure.
  \item Repeat the same v2 extraction and analysis for day 90 when the
        \texttt{spatial\_distribution\_v2\_\{c,d\}\_day90.mat} files are
        available.
\end{enumerate}

\end{document}
